\documentclass[12pt]{article}

\input{~/School/Notes/header.tex}

\title{Домашнее задание по алгебре}
\author{Денисов Илья}
\date{}

\begin{document}
	\maketitle

	\section*{Задача 1}

	\noindent
	\textbf{Условие.} Напишите многочлен с целыми коэффициентами, корнем которого является $a + \dfrac{1}{a}$, если $a^3 = a + 1$.

	\medskip
	\noindent
	\textbf{Идея.} Многочлен $f(x) = x^3 - x - 1$ унитарный, и $f(a) = 0$. Обозначим $c = a + \dfrac{1}{a}$. Сначала запишем $c$ как многочлен от $a$, затем применим метод матрицы умножения: $c$ является корнем характеристического многочлена матрицы умножения на $c$ в базисе $1, a, a^2$.

	\medskip
	\noindent
	\textbf{Шаг 1: выражаем $c$ через $a$.} Так как $f(0) = -1 \ne 0$, число $a$ не равно нулю. Из $a^3 = a + 1$ получаем $a(a^2 - 1) = 1$, откуда
	\[
		\frac{1}{a} = a^2 - 1, \qquad c = a^2 + a - 1.
	\]

	\medskip
	\noindent
	\textbf{Шаг 2: умножаем $c$ на базисные элементы.} Воспользуемся тем, что $a^3 = a + 1$ и $a^4 = a \cdot a^3 = a^2 + a$:
	\begin{align*}
		c \cdot 1 &= -1 + a + a^2, \\
		c \cdot a &= a^3 + a^2 - a = 1 + a^2, \\
		c \cdot a^2 &= a^4 + a^3 - a^2 = (a^2 + a) + (a + 1) - a^2 = 1 + 2a.
	\end{align*}
	Выпишем матрицу $M$: в $j$-й строке стоят коэффициенты при $1, a, a^2$ в разложении $c \cdot a^j$:
	\[
		M = \begin{pmatrix}
			-1 & 1 & 1 \\
			1 & 0 & 1 \\
			1 & 2 & 0
		\end{pmatrix}.
	\]

	\medskip
	\noindent
	\textbf{Шаг 3: характеристический многочлен.} Пусть $w = (1, a, a^2)^T \ne 0$. Равенства из шага 2 означают, что $Mw = c\,w$, то есть $(cE - M)w = 0$. Так как $w \ne 0$, получаем $\det(cE - M) = 0$: число $c$ является корнем многочлена $\chi_M(x) = \det(xE - M)$. Вычислим его, раскрывая определитель по первой строке:
	\begin{align*}
		\chi_M(x) &= \begin{vmatrix}
			x + 1 & -1 & -1 \\
			-1 & x & -1 \\
			-1 & -2 & x
		\end{vmatrix} \\
		&= (x + 1)(x^2 - 2) + (-x - 1) - (x + 2) \\
		&= x^3 + x^2 - 2x - 2 - x - 1 - x - 2 = x^3 + x^2 - 4x - 5.
	\end{align*}

	\medskip
	\noindent
	\textbf{Проверка.} Подставим $c = a^2 + a - 1$ непосредственно. Используя $a^3 = a + 1$ и $a^4 = a^2 + a$, получаем
	\begin{align*}
		c^2 &= a^4 + 2a^3 - a^2 - 2a + 1 = (a^2 + a) + 2(a + 1) - a^2 - 2a + 1 = a + 3, \\
		c^3 &= c \cdot c^2 = (a^2 + a - 1)(a + 3) = a^3 + 4a^2 + 2a - 3 = 4a^2 + 3a - 2.
	\end{align*}
	Тогда
	\[
		c^3 + c^2 - 4c - 5 = (4a^2 + 3a - 2) + (a + 3) - 4(a^2 + a - 1) - 5 = 0.
	\]

	\medskip
	\noindent
	\textbf{Замечание.} Число $a$ можно брать любым корнем $x^3 - x - 1$ (в том числе комплексным): рассуждение использует только равенство $a^3 = a + 1$. У найденного многочлена нет рациональных корней (кандидаты $\pm 1, \pm 5$ не подходят), поэтому как многочлен третьей степени он неприводим над $\mathbb{Q}$ и является минимальным многочленом числа $c$.

	\medskip
	\noindent
	\textbf{Ответ.}
	\[
		\boxed{x^3 + x^2 - 4x - 5}
	\]

	%==================================================
	\section*{Задача 2}

	\noindent
	\textbf{Условие.} Многочлен 5-й степени с рациональными коэффициентами не имеет рациональных корней. Докажите, что он не имеет кратных корней.

	\begin{proof}
		Пусть $f \in \mathbb{Q}[x]$, $\deg f = 5$, и у $f$ нет рациональных корней. Предположим противное: у $f$ есть кратный корень $\alpha \in \mathbb{C}$, то есть $(x - \alpha)^2 \mid f$ в $\mathbb{C}[x]$. Число $\alpha$ алгебраично над $\mathbb{Q}$, так как $f(\alpha) = 0$. Пусть $m \in \mathbb{Q}[x]$~--- его минимальный многочлен. Напомним, что $m$ неприводим над $\mathbb{Q}$, а неприводимый над $\mathbb{Q}$ многочлен не имеет кратных корней.

		\medskip
		\noindent
		\textbf{Шаг 1: $m$ делит $f$.} Так как $f(\alpha) = 0$, по теореме о минимальном многочлене $m \mid f$, то есть $f = mh$ при некотором $h \in \mathbb{Q}[x]$.

		\medskip
		\noindent
		\textbf{Шаг 2: $\alpha$ является корнем $h$.} По теореме Безу $m = (x - \alpha)u$ для некоторого $u \in \mathbb{C}[x]$. Если бы $u(\alpha) = 0$, то $(x - \alpha) \mid u$, и $(x - \alpha)^2 \mid m$, то есть $\alpha$ был бы кратным корнем неприводимого многочлена $m$. Это невозможно, значит, $u(\alpha) \ne 0$.

		Теперь $f = (x - \alpha)\,uh$ и $(x - \alpha)^2 \mid f$. Сократив на $x - \alpha$, получаем $(x - \alpha) \mid uh$, то есть $u(\alpha)h(\alpha) = 0$. Так как $u(\alpha) \ne 0$, имеем $h(\alpha) = 0$.

		\medskip
		\noindent
		\textbf{Шаг 3: $m^2$ делит $f$.} Многочлен $h \in \mathbb{Q}[x]$ обращается в $\alpha$ в ноль, поэтому по теореме о минимальном многочлене $m \mid h$. Значит, $h = mq$ и
		\[
			f = m^2 q, \qquad q \in \mathbb{Q}[x].
		\]

		\medskip
		\noindent
		\textbf{Шаг 4: подсчёт степеней.} Если $\deg m = 1$, то $m = x + b$ при $b \in \mathbb{Q}$ (многочлен унитарный), и $-b \in \mathbb{Q}$~--- рациональный корень $f$, что противоречит условию. Поэтому $\deg m \geqslant 2$. Из $m^2 \mid f$ следует $2\deg m \leqslant 5$, откуда $\deg m = 2$. Тогда
		\[
			\deg q = 5 - 2 \cdot 2 = 1.
		\]
		Многочлен $q = ax + b$ с $a, b \in \mathbb{Q}$, $a \ne 0$, имеет рациональный корень $r = -b/a$, и $f(r) = m(r)^2 q(r) = 0$. Мы получили рациональный корень $f$, что противоречит условию.

		Следовательно, у $f$ нет кратных корней.
	\end{proof}

	%==================================================
	\section*{Задача 3}

	\noindent
	\textbf{Условие.} Хулиган Петя украл из некоторой конечной группы $G$ треть её элементов и с постыдным удовлетворением отметил, что его добыча образует подгруппу. Докажите, что хулиган Вася может украсть у Пети не менее шестой части этой добычи, оформив у себя нормальную подгруппу $G$.

	\medskip
	\noindent
	\textbf{Формализация.} Пусть $H \leqslant G$~--- добыча Пети, $|H| = |G|/3$. Нужно найти подгруппу $N \trianglelefteq G$, нормальную в $G$, с $N \leqslant H$ и $|N| \geqslant |H|/6$. Мы докажем более сильную оценку:
	\[
		|N| \geqslant \frac{|G|}{6} = \frac{|H|}{2}.
	\]

	\begin{proof}
		\textbf{Шаг 1: индекс.} По теореме Лагранжа $|G| = |H| \cdot [G : H]$, поэтому $[G : H] = 3$. Значит, множество левых смежных классов $X = G/H$ состоит из трёх элементов.

		\medskip
		\noindent
		\textbf{Шаг 2: действие на смежных классах.} Для $g \in G$ определим отображение
		\[
			\tau_g: X \to X, \qquad \tau_g(xH) = gxH.
		\]
		Оно корректно: если $xH = yH$, то $y = xh$ при некотором $h \in H$, и $gyH = gxhH = gxH$. Далее, $\tau_e = \mathrm{id}$ и $\tau_g \circ \tau_{g'} = \tau_{gg'}$, откуда $\tau_g \circ \tau_{g^{-1}} = \tau_{g^{-1}} \circ \tau_g = \mathrm{id}$, то есть каждое $\tau_g$~--- перестановка множества $X$. Итак, отображение
		\[
			\varphi: G \to S_X, \qquad \varphi(g) = \tau_g,
		\]
		является гомоморфизмом групп, а в группе $S_X$ ровно $3! = 6$ элементов.

		\medskip
		\noindent
		\textbf{Шаг 3: ядро нормально.} Положим $N = \operatorname{Ker}\varphi = \{ g \in G \mid \tau_g = \mathrm{id} \}$. Это подгруппа $G$, а если $n \in N$ и $g \in G$, то
		\[
			\varphi(gng^{-1}) = \varphi(g)\,\varphi(n)\,\varphi(g)^{-1} = \varphi(g)\,\mathrm{id}\,\varphi(g)^{-1} = \mathrm{id},
		\]
		то есть $gng^{-1} \in N$. Значит, $N \trianglelefteq G$.

		\medskip
		\noindent
		\textbf{Шаг 4: $N$ лежит в $H$.} Пусть $g \in N$. Тогда $\tau_g(eH) = eH$, то есть $gH = H$, и $g = ge \in gH = H$. Итак, $N \leqslant H$: Вася действительно «крадёт» у Пети.

		\medskip
		\noindent
		\textbf{Шаг 5: оценка размера.} Для $g, h \in G$
		\[
			\varphi(g) = \varphi(h) \iff \varphi(g^{-1}h) = \mathrm{id} \iff g^{-1}h \in N \iff gN = hN.
		\]
		Поэтому $gN \mapsto \varphi(g)$~--- биекция между $G/N$ и образом $\operatorname{Im}\varphi$, то есть $[G : N] = |\operatorname{Im}\varphi|$. По теореме Лагранжа $|G| = |N| \cdot [G : N]$. Так как $\operatorname{Im}\varphi \subset S_X$, имеем $|\operatorname{Im}\varphi| \leqslant 6$, и
		\[
			|N| = \frac{|G|}{|\operatorname{Im}\varphi|} \geqslant \frac{|G|}{6} = \frac{3|H|}{6} = \frac{|H|}{2} \geqslant \frac{|H|}{6}.
		\]

		Таким образом, Вася может взять нормальную в $G$ подгруппу $N \leqslant H$, которая составляет не менее половины добычи Пети, а значит, и не менее её шестой части.
	\end{proof}
\end{document}